% ============================================================================
% Fichier  : partie4/P04_C17_cas_integrateurs.tex
% Titre    : Études de Cas Intégrateurs
% Auteur   : MINKA MI NGUIDJOÏ Thierry Emmanuel
% Version  : 1.0 -- octobre 2026
% ============================================================================
% Chapitre-cas. Chaque cas suit le plan : brief initial, investigation
% guidée, révélations progressives, synthèse et rapport. Le chapitre se
% termine par une grille d'erreurs et une simulation d'examen.
% Laboratoire associé : GCI-LAB-10 (diagnostic et feuille de route chiffrée).
% Tous les cas sont fictifs (noms, dates, montants) ; les mécanismes sont
% exacts.
% ============================================================================

\chapter{Études de Cas Intégrateurs}
\label{chap:cas-integrateurs}
\chaptermark{Études de cas intégrateurs}

\epigraph{%
    « On ne comprend vraiment un référentiel que le jour où il faut
    l'appliquer à une organisation qui ne lui ressemble pas. »%
}{Principe de l'apprentissage par les cas}

\niveauChapitre{Expert}{%
    Maîtrise des Chapitres~\ref{chap:cobit-capacite}
    (capacité et feuille de route), \ref{chap:audit} (audit et constats) et
    \ref{chap:processus} (processus, RACI, indicateurs).%
}

\begin{boxobjectifs}
\begin{itemize}
    \item Mobiliser l'ensemble du cours sur les quatre cas fil rouge.
    \item Produire un diagnostic de gouvernance complet.
    \item Bâtir une feuille de route COBIT--ITIL chiffrée.
    \item Présenter et défendre ses recommandations.
    \item Se préparer à l'examen (grille d'erreurs, simulation).
\end{itemize}
\end{boxobjectifs}

\bigskip

% ============================================================================
\section*{Mode d'emploi du chapitre}
% ============================================================================

Ce chapitre ne présente aucune notion nouvelle. Il reprend les quatre
organisations du cours, un an après les événements qui les ont fait
connaître, et pose à chaque fois la même question~: \emph{la gouvernance
des SI a-t-elle réellement progressé, et que faut-il faire maintenant~?}

Chaque cas suit quatre temps. Le \textbf{brief} donne la situation et la
demande. L'\textbf{investigation guidée} propose les questions à se poser,
avec le chapitre qui permet d'y répondre. Les \textbf{révélations} apportent,
par étapes, des faits nouveaux qui obligent à réviser le premier
jugement. La \textbf{synthèse} donne une réponse de référence, à comparer
avec la vôtre. Travaillez chaque cas \emph{avant} de lire sa synthèse.

La section~\ref{sec:c17-methode} fixe la méthode commune. Les
sections~\ref{sec:c17-sanaga} à~\ref{sec:c17-noun} traitent les quatre
cas. Les sections~\ref{sec:c17-erreurs} et~\ref{sec:c17-examen}
proposent une grille d'erreurs et une simulation d'examen. Le laboratoire
GCI-LAB-10 de NDAP SUKULU reprend la démarche sur les quatre dossiers.

\separateur

% ============================================================================
\section{La Méthode du Diagnostic Intégré}
\label{sec:c17-methode}
% ============================================================================

\subsection{Six volets}

Un diagnostic de gouvernance des SI couvre six volets. Chacun mobilise une
partie du cours.

\begin{tableaucours}{L{3.0cm}L{5.2cm}L{4.0cm}}{%
    \tableauHeader{Volet} &
    \tableauHeader{Questions} &
    \tableauHeader{Chapitres}
}
1. Organes et décisions & Qui décide de quoi~? Les organes fonctionnent-ils~?
& \ref{chap:gouv-entreprise}, \ref{chap:alignement} \\
2. Valeur et portefeuille & Les investissements sont-ils arbitrés et leurs
bénéfices suivis~? & \ref{chap:alignement} \\
3. Système COBIT & Objectifs prioritaires, capacité constatée et cible,
écarts & \ref{chap:cobit-objectifs} à~\ref{chap:cobit-capacite} \\
4. Conformité & Obligations légales et réglementaires, preuves &
\ref{chap:normes}, \ref{chap:droit} \\
5. Services & Pratiques ITIL, niveaux de service, indicateurs &
\ref{chap:itil-svs} à~\ref{chap:processus} \\
6. Fournisseurs et continuité & Contrats, supervision, plans testés &
\ref{chap:itil-niveaux} \\
\end{tableaucours}
\vspace{-0.8em}
\begin{center}{\small\itshape Tableau~17.1. Les six volets du diagnostic
intégré}\end{center}

\subsection{Le livrable}

\begin{boxprocedure}
\textbf{Le rapport de diagnostic.} Le rapport remis à l'organe de
gouvernance comporte~:
\begin{itemize}
    \item une synthèse d'une page, avec trois messages et trois décisions
          demandées~;
    \item le périmètre, les sources et la méthode (entretiens, pièces,
          tests)~;
    \item pour chaque volet, les constats en cinq éléments
          (Chapitre~\ref{chap:audit}), chacun avec son niveau de
          confiance~;
    \item le tableau des écarts de capacité (Chapitre~\ref{chap:cobit-capacite})~;
    \item la feuille de route chiffrée, avec ses propriétaires et ses
          échéances~;
    \item les risques de la feuille de route elle-même.
\end{itemize}
\end{boxprocedure}

% ============================================================================
\section{Cas 1 : la BANQUE SANAGA, un An Après}
\label{sec:c17-sanaga}
% ============================================================================

\subsection{Brief initial}

\begin{boxscenario}{Opération SANAGA. Mars 2027}
Un an après la panne, la COBAC annonce une mission de contrôle sur place
consacrée au risque informatique. Le président du conseil
d'administration vous demande, comme consultant indépendant, un
diagnostic préalable~: « Avons-nous fait ce que nous avions annoncé à la
COBAC, et sommes-nous prêts~? »

\medskip
\textbf{Pièces remises.} Le rapport circonstancié à la COBAC (avril~2026),
la feuille de route à dix-huit mois (Chapitre~\ref{chap:cobit-capacite}),
les procès-verbaux du comité stratégique des SI, le registre des
changements de l'année, le contrat renégocié avec DIGICORE AFRIQUE, le
compte rendu de l'agence pilote de la nouvelle bascule.
\end{boxscenario}

\subsection{Investigation guidée}

\begin{enumerate}
    \item Le comité stratégique des SI fonctionne-t-il selon les critères du
          Chapitre~\ref{chap:gouv-entreprise}~? Reçoit-il une information
          indépendante de celle du directeur général~?
    \item Où en est chaque ligne de la feuille de route~? Quelles capacités
          ont été réévaluées, et par qui (Chapitre~\ref{chap:cobit-capacite})~?
    \item Le registre des changements montre-t-il une gestion conforme à la
          grille de risque (Chapitre~\ref{chap:itil-changements})~?
    \item Le contrat de DIGICORE AFRIQUE contient-il les clauses essentielles
          (Chapitre~\ref{chap:itil-niveaux}), et sont-elles appliquées~?
    \item Les obligations légales et réglementaires sont-elles tenues et
          prouvées (Chapitre~\ref{chap:droit})~?
\end{enumerate}

\subsection{Révélations progressives}

\begin{boxscenario}{Opération SANAGA. Première révélation}
Le comité stratégique des SI s'est réuni quatre fois, comme prévu. Mais
ses procès-verbaux montrent qu'il n'a examiné que des présentations du
directeur général~; la DSI et l'audit interne n'ont jamais été entendus
séparément. Le pouvoir d'accès direct prévu par la charte n'a pas été
exercé.
\end{boxscenario}

\begin{boxscenario}{Opération SANAGA. Deuxième révélation}
Le registre compte 412~changements dans l'année, dont 91 qualifiés
d'urgents (22~\%). La procédure des changements urgents a bien été
rédigée, mais 64 de ces 91~changements n'ont pas fait l'objet de la revue
a posteriori prévue.
\end{boxscenario}

\begin{boxscenario}{Opération SANAGA. Troisième révélation}
Le contrat renégocié avec DIGICORE AFRIQUE contient désormais des
engagements de service, des pénalités et un droit d'audit. Les pénalités
dues pour deux manquements en 2026 (évaluées à 38~millions FCFA) n'ont pas
été réclamées~: la direction des achats « ne voulait pas dégrader la
relation pendant la nouvelle bascule ». Le droit d'audit n'a jamais été
exercé.
\end{boxscenario}

\subsection{Synthèse et rapport}

\begin{tableaucours}{L{3.2cm}L{5.2cm}cL{2.4cm}}{%
    \tableauHeader{Volet} &
    \tableauHeader{Constat principal} &
    \tableauHeader{Conf.} &
    \tableauHeader{Objectif}
}
Organes & Le comité existe, mais reproduit l'asymétrie d'information
qu'il devait réduire & Élevée & EDM01 \\
COBIT & Gains rapides réalisés (DSS04, MEA02)~; chantiers de fond en retard
(EDM03, APO12) & Moyenne & EDM03, APO12 \\
Changements & Taux d'urgence élevé, revues a posteriori non faites (70~\%)
& Élevée & BAI06 \\
Fournisseurs & Clauses présentes mais non appliquées & Élevée & APO10 \\
Conformité & Audit de sécurité légal réalisé en 2026 & Élevée & MEA03 \\
\end{tableaucours}
\vspace{-0.8em}
\begin{center}{\small\itshape Tableau~17.2. Synthèse du diagnostic de la
BANQUE SANAGA}\end{center}

\textbf{Message principal au conseil.} La banque a mis en place les
\emph{structures} et les \emph{documents} annoncés, mais pas encore les
\emph{comportements}~: le comité n'écoute qu'une voix, les contrôles sont
contournés quand ils dérangent, les clauses ne sont pas appliquées. C'est
exactement ce que la grille de De Haes et Van Grembergen permet de
diagnostiquer (Chapitre~\ref{chap:gouv-entreprise})~: des structures sans
mécanismes relationnels. La COBAC le verra.

\begin{boxscenario}{Opération SANAGA. Feuille de route}
\begin{tableaucours}{L{5.6cm}L{2.4cm}rL{1.6cm}}{%
    \tableauHeader{Action} &
    \tableauHeader{Propriétaire} &
    \tableauHeader{Coût (M FCFA)} &
    \tableauHeader{Échéance}
}
Audition séparée de la DSI et de l'audit interne à chaque comité &
Président du comité & 0 & Immédiat \\
Revue a posteriori obligatoire des changements urgents, indicateur
trimestriel & Responsable des changements & 15 & 3 mois \\
Réclamation des pénalités 2026 et premier audit de DIGICORE AFRIQUE &
Direction des achats, RSSI & 25 & 3 mois \\
Appétence au risque approuvée par le conseil, registre des risques SI relié
au registre de la banque & Directeur des risques & 60 & 6 mois \\
Réévaluation indépendante de la capacité des dix objectifs & Audit interne
et cabinet externe & 45 & 9 mois \\
Formation des administrateurs à la gouvernance des SI & Secrétaire du
conseil & 12 & 6 mois \\
\textbf{Total} & & \textbf{157} & \\
\end{tableaucours}
\medskip
Le coût est à rapporter aux 1,9~milliard FCFA de la panne de 2026~: moins
de 9~\%.
\end{boxscenario}

% ============================================================================
\section{Cas 2 : KRIBILOG, la Certification en Jeu}
\label{sec:c17-kribi}
% ============================================================================

\subsection{Brief initial}

\begin{boxscenario}{Opération KRIBI. Mois 12 sur 18}
KRIBILOG SA doit obtenir la certification ISO/IEC~20000-1 dans six mois
pour conserver le contrat de l'armateur (Chapitre~\ref{chap:normes}).
M.~MBARGA vous demande un avis~: « Pouvons-nous solliciter l'organisme de
certification comme prévu au mois~16~? »

\medskip
\textbf{Pièces remises.} Le plan de certification, l'avenant au contrat de
NETPORT SERVICES, les indicateurs mensuels des six derniers mois, le
registre des problèmes, le rapport d'audit interne du mois~11.
\end{boxscenario}

\subsection{Investigation guidée}

\begin{enumerate}
    \item Les pratiques des articles~8.5 et~8.6 de la norme
          fonctionnent-elles depuis assez longtemps pour produire des
          enregistrements probants (Chapitre~\ref{chap:normes})~?
    \item Les indicateurs ont-ils progressé, et sont-ils mesurés par
          KRIBILOG elle-même (Chapitres~\ref{chap:itil-incidents}
          et~\ref{chap:itil-niveaux})~?
    \item Le correctif de l'éditeur a-t-il été installé, et le problème
          clos (Chapitre~\ref{chap:itil-changements})~?
    \item L'avenant donne-t-il à KRIBILOG la maîtrise d'un service exploité
          par un tiers~?
\end{enumerate}

\subsection{Révélations progressives}

\begin{boxscenario}{Opération KRIBI. Première révélation}
Les progrès sont réels. Le taux de réouverture est passé de 30~\% à 12~\%,
le volume mensuel de tickets de 4~800 à 3~100, grâce au libre-service et à
l'installation du correctif. La disponibilité de bout en bout du service
de suivi atteint 99,3~\% en moyenne sur six mois, avec un second
opérateur mobile depuis le mois~8.
\end{boxscenario}

\begin{boxscenario}{Opération KRIBI. Deuxième révélation}
L'audit interne du mois~11 relève que NETPORT SERVICES a confié, sans en
informer KRIBILOG, l'hébergement des sauvegardes à un sous-traitant situé
à Casablanca. L'avenant ne traite ni de la sous-traitance en cascade, ni de
la localisation des données.
\end{boxscenario}

\begin{boxscenario}{Opération KRIBI. Troisième révélation}
La revue de direction prévue au mois~15 par la norme n'est pas encore
programmée~; le directeur général « n'en voit pas l'intérêt, puisque les
indicateurs sont bons ».
\end{boxscenario}

\subsection{Synthèse et rapport}

\textbf{Avis.} Le calendrier peut être tenu, \emph{sous deux conditions}. La
première tient au fournisseur~: un nouvel avenant doit traiter la
sous-traitance et la localisation des données avant l'audit, faute de
quoi l'organisme de certification relèvera une non-conformité majeure sur
la maîtrise des processus confiés à des tiers, et KRIBILOG s'exposera au
titre de la loi n°~2024/017 (Chapitre~\ref{chap:droit}). La seconde tient à
la direction~: la revue de direction est une exigence de la norme
(article~9 de la structure harmonisée), et non une formalité~; sans elle,
pas de certification.

\begin{boxmetier}{Responsable des services IT}
\textbf{Présenter l'avis au directeur général}
\begin{enumerate}
    \item Commencer par les progrès, chiffrés~: ils sont réels et
          mesurés.
    \item Présenter les deux conditions comme des risques pour le contrat
          de l'armateur, en termes d'affaires et non de norme.
    \item Proposer la date de la revue de direction et son ordre du jour.
    \item Demander une décision écrite sur l'avenant avec NETPORT
          SERVICES.
\end{enumerate}
\end{boxmetier}

% ============================================================================
\section{Cas 3 : l'ANPP face à la Chambre des Comptes}
\label{sec:c17-mefou}
% ============================================================================

\subsection{Brief initial}

\begin{boxscenario}{Opération MEFOU. Le rapport provisoire}
La Chambre des comptes a transmis à l'ANPP son rapport provisoire d'audit
des SI (Chapitre~\ref{chap:audit}). Le directeur général dispose de deux
mois pour répondre. Il vous demande de préparer la réponse et le plan
d'action.

\medskip
\textbf{Constats du rapport provisoire.} Absence de comité stratégique des
SI jusqu'en 2026~; projets lancés sans dossier d'investissement~; trois
projets ayant consommé 71~\% du budget du schéma directeur sans être
achevés~; 14~comptes d'agents partis actifs sur le système de paiement~;
séparation des tâches défaillante sur 6~comptes sur 25.
\end{boxscenario}

\subsection{Investigation guidée}

\begin{enumerate}
    \item Quels constats sont fondés, et lesquels peuvent être nuancés par
          des pièces que l'auditeur n'a pas vues (Chapitre~\ref{chap:audit})~?
    \item Comment réarbitrer le portefeuille restant (Chapitre~\ref{chap:alignement})~?
    \item Comment rattacher les projets aux programmes budgétaires
          (loi n°~2018/012)~?
    \item Quelles actions immédiates sur les accès (Chapitres~\ref{chap:itil-incidents}
          et~\ref{chap:audit})~?
\end{enumerate}

\subsection{Révélations progressives}

\begin{boxscenario}{Opération MEFOU. Première révélation}
Le projet de paie dématérialisée, présenté comme « inachevé » dans le
rapport, est en réalité en service depuis janvier~2027. Mais aucun suivi
des bénéfices n'a été fait~: personne ne sait si les délais de paiement
ont diminué.
\end{boxscenario}

\begin{boxscenario}{Opération MEFOU. Deuxième révélation}
Les 14~comptes actifs d'agents partis ont été désactivés le lendemain du
passage de l'auditeur. Mais la cause (l'absence de notification des
départs par les ressources humaines) n'a pas été traitée~: trois nouveaux
départs, depuis, n'ont pas été signalés.
\end{boxscenario}

\subsection{Synthèse et rapport}

\textbf{La réponse à la Chambre des comptes} admet les constats fondés,
apporte les pièces qui nuancent le constat sur le projet de paie, et
présente un plan d'action qui traite les \emph{causes}~:
\begin{itemize}
    \item un comité stratégique des SI, créé par décision du conseil, avec
          une charte inspirée de l'ISO/IEC~38500 (Chapitre~\ref{chap:normes})~;
    \item un réarbitrage du portefeuille restant par score multicritère
          (Chapitre~\ref{chap:alignement}), et l'arrêt d'un projet~;
    \item un suivi des bénéfices du projet de paie, avec un indicateur
          rattaché au projet de performance annuel du programme~;
    \item une notification automatique des départs et une revue trimestrielle
          des accès (DSS05, DSS06).
\end{itemize}

\begin{boxalert}
\textbf{Violation typique.} Corriger le symptôme le lendemain de l'audit
(désactiver les 14~comptes) et présenter la correction comme une réponse
au constat. L'auditeur vérifiera la cause, et constatera la récidive.

\smallskip
\textbf{Règle~:} une réponse à un rapport d'audit traite la cause, désigne
un responsable et une échéance, et indique la pièce qui prouvera la mise
en œuvre.
\end{boxalert}

% ============================================================================
\section{Cas 4 : NOUN PAY, Croître sans Perdre la Confiance}
\label{sec:c17-noun}
% ============================================================================

\subsection{Brief initial}

\begin{boxscenario}{Opération NOUN. La levée de fonds}
NOUN PAY prépare une levée de fonds de 2~milliards FCFA. Les investisseurs
exigent un diagnostic de gouvernance des SI avant de signer. Au même
moment, l'Autorité de protection des données poursuit son enquête sur le
modèle de notation (Chapitre~\ref{chap:droit}).

\medskip
\textbf{Pièces remises.} La matrice RACI du processus de mise à jour du
modèle (Chapitre~\ref{chap:processus}), le contrat type du fournisseur de
nuage, la réponse provisoire à l'Autorité, les statistiques de refus de
crédit par région.
\end{boxscenario}

\subsection{Investigation guidée}

\begin{enumerate}
    \item Quels facteurs de conception caractérisent NOUN PAY, et quel
          système de gouvernance proportionné en découle
          (Chapitre~\ref{chap:cobit-conception})~?
    \item Le contrat du fournisseur de nuage permet-il de maîtriser la
          localisation des données, la réversibilité et l'audit
          (Chapitre~\ref{chap:itil-niveaux})~?
    \item Qui répond, dans l'entreprise, des effets du modèle sur les
          clients (Chapitre~\ref{chap:processus})~?
\end{enumerate}

\subsection{Révélations progressives}

\begin{boxscenario}{Opération NOUN. Première révélation}
Les statistiques montrent que le taux de refus est de 61~\% pour les
demandeurs d'une région, contre 24~\% en moyenne nationale. Le modèle
utilise, parmi ses variables, l'opérateur mobile du demandeur, très
corrélé à la région.
\end{boxscenario}

\begin{boxscenario}{Opération NOUN. Deuxième révélation}
Le contrat type du fournisseur de nuage ne permet ni de choisir la
localisation des données, ni d'auditer le fournisseur~; il renvoie à ses
propres certifications. Un contrat d'entreprise négocié existe chez le
même fournisseur, pour un surcoût annuel de 18~millions FCFA.
\end{boxscenario}

\subsection{Synthèse et rapport}

\textbf{Message aux investisseurs.} NOUN PAY n'a pas besoin d'un système
COBIT complet, mais de six objectifs bien tenus (Chapitre~\ref{chap:cobit-conception})~:
EDM03 (appétence au risque), APO10 (fournisseur de nuage), APO12
(risques), APO14 (données), BAI06 (changements du modèle) et MEA03
(conformité). Deux décisions sont urgentes~: souscrire le contrat
d'entreprise du fournisseur de nuage, et suspendre l'usage de la variable
« opérateur mobile » en attendant une analyse d'impact du modèle. Le
Chapitre~\ref{chap:donnees-ia} reprendra ce cas sous l'angle de la
gouvernance de l'intelligence artificielle.

% ============================================================================
\section{Grille des Erreurs Transversales}
\label{sec:c17-erreurs}
% ============================================================================

Les quatre cas font réapparaître, sous des formes différentes, les mêmes
erreurs. La grille du Chapitre~\ref{chap:defaillances} s'enrichit de
celles que révèle la mise en œuvre.

\begin{tableaucours}{L{4.4cm}L{2.4cm}L{5.4cm}}{%
    \tableauHeader{Erreur} &
    \tableauHeader{Cas} &
    \tableauHeader{Antidote}
}
Des structures sans comportements & SANAGA & Évaluer le fonctionnement des
organes, pas leur existence (Chapitre~\ref{chap:cobit-capacite}) \\
Des contrôles contournés quand ils gênent & SANAGA & Indicateur de
contournement (taux d'urgence, revues non faites) \\
Des clauses non appliquées & SANAGA & Comité de service, réclamation
systématique des pénalités \\
Une sous-traitance en cascade ignorée & KRIBI & Clause de sous-traitance et
de localisation, droit d'audit \\
Un indicateur bon qui dispense de la revue & KRIBI & La revue de
direction est une exigence, pas une option \\
Un symptôme corrigé, une cause intacte & MEFOU & Traiter la cause, prouver
par une pièce \\
Un projet livré sans suivi des bénéfices & MEFOU & EDM02, indicateur
rattaché au programme \\
Une décision automatisée sans responsable & NOUN & Un A pour l'évaluation
de l'impact sur les clients \\
\end{tableaucours}
\vspace{-0.8em}
\begin{center}{\small\itshape Tableau~17.3. Erreurs transversales de mise en
œuvre}\end{center}

% ============================================================================
\section{Simulation d'Examen}
\label{sec:c17-examen}
% ============================================================================

\begin{boxinfo}
\textbf{Sujet type.} \emph{Durée~: 3~heures. Documents autorisés~: aucun,
sauf le tableau des 40~objectifs COBIT fourni.}

\medskip
\textbf{Cas} \textit{(fictif, construit pour l'examen)}. La SOCIÉTÉ D'EAU
DU NORD, entreprise publique
installée à Garoua, distribue l'eau potable à 180~000~abonnés. Son système
de facturation, exploité par un prestataire, a émis en septembre des
factures erronées à 30~000~abonnés. Le conseil d'administration, qui se
réunit deux fois par an, l'a appris par la presse. La direction
informatique compte huit personnes. Aucun comité ne traite du numérique.

\medskip
\begin{enumerate}
    \item \textbf{(4~points)} Distinguez la cause technique probable et les
          causes de gouvernance, en appliquant les trois questions réflexes.
    \item \textbf{(4~points)} Proposez les organes de gouvernance des SI
          adaptés à une entreprise publique de cette taille, en tenant
          compte de la tutelle.
    \item \textbf{(4~points)} Renseignez les facteurs de conception et
          retenez six objectifs COBIT prioritaires, avec leur capacité
          cible.
    \item \textbf{(4~points)} Rédigez deux constats d'audit en cinq éléments,
          l'un sur la gestion des changements, l'autre sur le contrat du
          prestataire.
    \item \textbf{(4~points)} Présentez une feuille de route à douze mois,
          chiffrée en FCFA, en distinguant gains rapides et chantiers de
          fond.
\end{enumerate}
\end{boxinfo}

\begin{boxalert}
\textbf{Ce que le correcteur attend.} Des constatations séparées des
analyses~; des références précises (objectifs COBIT, pratiques ITIL,
textes de loi vérifiés)~; des recommandations qui traitent des causes~; une
gouvernance \emph{proportionnée} à une organisation de huit informaticiens.
Une copie qui recommande de « déployer COBIT et ITIL » sans choisir
n'obtient pas la moyenne.
\end{boxalert}

% ============================================================================
\separateur
% ============================================================================

\begin{boxresume}
\textbf{Ce qu'il faut retenir de ce chapitre~:}
\begin{itemize}
    \item \textbf{Six volets}~: organes et décisions, valeur, système COBIT,
          conformité, services, fournisseurs et continuité. Un diagnostic
          les couvre tous.
    \item \textbf{Les structures ne suffisent pas}~: un an après, la BANQUE
          SANAGA a ses comités et ses documents, mais pas encore les
          comportements.
    \item \textbf{Les bons indicateurs ne dispensent de rien}~: KRIBILOG
          progresse, mais une sous-traitance cachée et une revue de
          direction négligée menacent sa certification.
    \item \textbf{Traiter les causes}~: l'ANPP corrige des symptômes le
          lendemain de l'audit et voit les écarts revenir.
    \item \textbf{La proportion}~: NOUN PAY n'a besoin que de six objectifs
          bien tenus, et de deux décisions urgentes.
    \item \textbf{Une feuille de route se chiffre}~: à rapporter au coût des
          défaillances qu'elle prévient.
\end{itemize}
\end{boxresume}

\bigskip

\begin{boxexo}[title={\ding{45}\quad Exercice~17.1. Le rapport de la BANQUE SANAGA \niveaubase}]
Rédigez la synthèse d'une page du diagnostic de la BANQUE SANAGA~: trois
messages, trois décisions demandées au conseil, et le coût de la feuille
de route rapporté au coût de la panne.
\end{boxexo}

\begin{boxexo}[title={\ding{45}\quad Exercice~17.2. L'avenant de KRIBILOG \niveaumoyen}]
Rédigez les clauses de l'avenant au contrat de NETPORT SERVICES qui
traitent la sous-traitance en cascade, la localisation des sauvegardes et
le droit d'audit. Pour chacune, indiquez le risque traité et la pièce qui
permettra de vérifier son application.
\end{boxexo}

\begin{boxexo}[title={\ding{45}\quad Exercice~17.3. Traiter le sujet d'examen \niveauexpert}]
Traitez en trois heures le sujet de la section~\ref{sec:c17-examen}, puis
comparez votre copie avec les attentes du correcteur. Pour chaque
question, indiquez les chapitres du cours que vous avez mobilisés.
\end{boxexo}

\bigskip

\begin{boxinfo}
\textbf{Transition.} La Partie~\ref{part:pratique} s'achève. Les quatre cas
ont montré que la gouvernance des SI se joue sur des terrains qui
évoluent vite~: l'informatique en nuage, les méthodes agiles, les données
et l'intelligence artificielle. La Partie~\ref{part:avance} les aborde,
en commençant par le Chapitre~\ref{chap:agile-cloud}~: gouverner à grande
vitesse.
\end{boxinfo}
